% Independently written from Illusie I, 2.2.1--2.2.10, printed pp.25--38.
% Local labels only. Source snippet for bounded root-chapter composition.

\section{Free abelian sheaves and $n$-equivalences}
\label{section-illusie-I-whitehead-free-abelian}

\begin{definition}
\label{definition-illusie-I-n-equivalence}
Let $\mathcal{T}$ be a topos and let $f : X \to Y$ be a map of
simplicial objects of $\mathcal{T}$. For $n \geq 0$ we say that $f$ is an
{\it $n$-equivalence} if $\pi_0(f)$ is an epimorphism and, after restriction
to every object $U$ and choice of a base point $x$ of $X|_U$, the map
$$
\pi_i(X|_U, x) \longrightarrow \pi_i(Y|_U, f(x))
$$
is an isomorphism for $i < n$ and an epimorphism for $i = n$.
We say that $f$ is an {\it equivalence} if it is an $n$-equivalence for
every $n$.
\end{definition}

\begin{lemma}[Relative Hurewicz for free simplicial abelian groups]
\label{lemma-illusie-I-relative-hurewicz-free-abelian}
Let $n \geq 2$ and let $f : K \to L$ be an $n$-equivalence of simplicial
sets. Then
$$
\mathbf{Z}[f] : \mathbf{Z}[K] \longrightarrow \mathbf{Z}[L]
$$
is an $n$-equivalence of simplicial abelian groups.
\end{lemma}

\begin{proof}
The assertion can be checked one connected component at a time. Replace the
geometric realization of $f$ by the inclusion of its mapping cylinder. The
$n$-equivalence condition and the long exact sequence of relative homotopy
groups say that the resulting pair has zero relative homotopy groups through
degree $n$. Since $n \geq 2$, the relative Hurewicz theorem applies on every
component and gives
$$
H_i(|L|, |K|; \mathbf{Z}) = 0 \quad\text{for }i \leq n.
$$
The normalized chain complex of $\mathbf{Z}[K]$ is the simplicial chain
complex of $K$, and similarly for $L$. Hence the displayed vanishing says
that its induced map on homology is an isomorphism below $n$ and a
surjection in degree $n$. Homotopy groups of a simplicial abelian group are
the homology groups of its normalized complex. Finally, the
$n$-equivalence condition makes $\pi_0(f)$ bijective, while
$\pi_0(\mathbf{Z}[K])$ is the free abelian group on $\pi_0(K)$. This proves
the lemma.
\end{proof}

\begin{lemma}[Local filtered colimits]
\label{lemma-illusie-I-local-filtered-colimit}
Let $\mathcal{T}$ be a topos. Let $L \to \mathcal{T}^{\mathrm{opp}}$ be a
cofibred category which is locally nonempty and in which any two objects
have a common successor locally and any two parallel arrows become equal
after a further local successor. Write $\operatorname{LSys}(L)$ for the
category of systems of sets indexed by $L$, compatible with arrows and
restriction in $\mathcal{T}$. Then the constant-system functor
$\mathcal{T} \to \operatorname{LSys}(L)$ has a left adjoint
$$
\colim_L : \operatorname{LSys}(L) \longrightarrow \mathcal{T}.
$$
This functor preserves all colimits and finite limits. A locally cofinal full
cofibred subcategory $L' \subset L$ gives the same local colimit.
Moreover, free abelian objects and the homotopy sheaves of pointed
simplicial objects commute with $\colim_L$.
\end{lemma}

\begin{proof}
Give $L^{\mathrm{opp}}$ the topology in which a family is covering exactly
when its image in $\mathcal{T}$ is covering. Local nonemptiness, common
successors, and local equality of parallel arrows are respectively the
density, local fullness, and local faithfulness conditions which identify
the sheaf topos of this site with $\mathcal{T}$. Under this identification,
$\colim_L$ is sheafification of the presheaf colimit of a local system. This
proves the adjunction and preservation of colimits. Filtered colimits of sets
commute with finite limits, and sheafification is left exact, so $\colim_L$
preserves finite limits as well.

The same site description for $L'$ gives the cofinality assertion: every
object is locally reached from $L'$, and the two local filteredness axioms
make the resulting identifications independent of the chosen arrows.

The free abelian object functor is a left adjoint, so its universal property
commutes with the displayed left adjoint. For homotopy sheaves, use the
construction in the proof of Lemma
\ref{lemma-illusie-I-homotopy-sheaves-colimits-pullback}: it uses finite
limits, colimits, and filtered colimits only. Each of these is preserved by
$\colim_L$, proving the last assertion.
\end{proof}

\begin{lemma}[Local finite-type approximation]
\label{lemma-illusie-I-local-finite-type-approximation}
Let $A$ be a ring, viewed as a constant ring object of a topos
$\mathcal{T}$, and let $u : E \to F$ be a map of complexes of
$A$-modules.
\begin{enumerate}
\item There are a locally filtered category $L$ and a local system
$u_i : E_i \to F_i$ of maps of complexes of finitely presented $A$-modules
such that $u = \colim_L u_i$.
\item If $E$ and $F$ vanish in degrees greater than $a$, the system can be
chosen with every $E_i$ and $F_i$ vanishing in degrees greater than $a$.
\item If $A$ is Noetherian, $E$ and $F$ vanish in degrees greater than
$a$, and
$$
H^q(\operatorname{Cone}(u)) = 0 \quad\text{for }q \geq b,
$$
then the full local subsystem of maps satisfying the same vanishing is
locally cofinal.
\end{enumerate}
\end{lemma}

\begin{proof}
Over an object $U$ of $\mathcal{T}$, consider maps $u_i \to u|_U$ where
$u_i$ is a map between complexes of finitely presented $A$-modules.
Restriction gives a cofibred category. A finite collection of sections of
$E$ and $F$, together with finitely many equations expressing compatibility
with the differentials and $u$, is represented locally by finite free
$A$-modules. Two such finite presentations have a common refinement, and
two maps between them which agree in $u$ agree after a further refinement.
Thus this category is locally filtered. Every section and every relation is
eventually represented, which proves that its local colimit is $u$. If
$E$ and $F$ vanish above $a$, discard the corresponding terms in every
finite presentation; this gives a locally cofinal subsystem and proves (2).

We record the two extension operations used for (3). Let
$v : C \to D$ be one finite presentation over an object $U$, with structure
maps $\alpha : C \to E|_U$ and $\beta : D \to F|_U$.

First suppose that $H^{q + 1}(u)$ is injective and that
$H^q(\operatorname{Cone}(v))$ is of finite type. The cohomology sequence of
the cone makes $\ker H^{q + 1}(v)$ finite type. Choose a finite free module
$M$ and cycles $c : M \to Z^{q + 1}(C)$ whose classes generate this kernel.
Choose $t : M \to D^q$ with $d t = v c$. Injectivity of
$H^{q + 1}(u)$ says that $\alpha c$ is locally a boundary; after a covering
choose $s : M \to E^q|_U$ with $d s = \alpha c$, and put
$\delta = u s - \beta t$. Thus $d\delta = 0$.

Adjoin $M$ in degree $q$ to both $C$ and $D$. On the source use outgoing
differential $(c,d)$, on the target use $(0,d)$, and in degree $q$ use the
map
$$
\begin{pmatrix}1&0\\ t&v\end{pmatrix} :
M \oplus C^q \longrightarrow M \oplus D^q.
$$
The structure maps in that degree are $(s,\alpha)$ and $(\delta,\beta)$.
These formulas define a refinement $v'$ of $v$. Its quotient by $v$ is the
identity map of the complex with $M$ in degree $q$ and zero elsewhere.
The cohomology sequence now shows that $H^{q + 1}(v')$ is injective and has
the same image as $H^{q + 1}(v)$.

Second suppose that $H^q(u)$ is surjective, $H^{q + 1}(v)$ is injective,
and $H^q(\operatorname{Cone}(v))$ is of finite type. The cone sequence
makes $H^q(D) \to H^q(\operatorname{Cone}(v))$ surjective. Choose a finite
free $M$ and $t : M \to Z^q(D)$ whose classes generate the latter module.
After a covering, surjectivity of $H^q(u)$ gives
$s : M \to Z^q(E|_U)$ lifting their images. Then
$\delta = u s - \beta t$ is a boundary; after another covering choose
$\epsilon : M \to F^{q - 1}|_U$ with $d\epsilon = \delta$.

Adjoin $M$ to $C^q$, and adjoin the two-term identity complex
$M \xrightarrow{1} M$ in degrees $q - 1,q$ to $D$. In degrees $q - 1,q$
the new map is
$$
C^{q - 1} \longrightarrow M \oplus D^{q - 1},\quad
x \longmapsto (0,vx),
$$
and
$$
M \oplus C^q \longrightarrow M \oplus D^q,\quad
(m,x) \longmapsto (m,t(m)+v(x)).
$$
Use $(s,\alpha)$ on the source in degree $q$, and
$(\epsilon,\beta)$ and $(\delta,\beta)$ on the target in degrees
$q - 1$ and $q$. These formulas define a refinement $v'$. Its quotient by
$v$ is the map from the complex with $M$ in degree $q$ to the identity
two-term complex just displayed. The cohomology sequence shows that
$H^{q + 1}(v')=H^{q + 1}(v)$ and that $H^q(v')$ is surjective.

Return to a finite presentation bounded above by $a$. Its cone cohomology
is of finite type because $A$ is Noetherian. Starting with the automatic
injectivity in degree $a + 1$, apply the second operation in degrees
$q=a,a-1,\ldots,b$. After the operation in degree $q>b$, apply the first
operation with index $q-1$; this makes $H^q(v)$ injective without changing
its already surjective image. At the end, $H^q(v)$ is an isomorphism for
$q>b$ and an epimorphism for $q=b$, equivalently its cone is acyclic in
degrees at least $b$. Thus every object of the original indexing category
locally maps to one in the asserted full subsystem. It is therefore locally
cofinal, proving (3).
\end{proof}

\begin{lemma}
\label{lemma-illusie-I-local-colimit-free-modules}
Let $L$ be locally filtered over $\mathcal{T}$. Let $(S_i)_{i \in L}$ and
$(A_i)_{i \in L}$ be local systems of sets and rings, and put
$$
S = \colim_L S_i, \qquad A = \colim_L A_i.
$$
Then the canonical map of $A$-modules
$$
\colim_L A_i^{(S_i)} \longrightarrow A^{(S)}
$$
is an isomorphism.
\end{lemma}

\begin{proof}
A local section of $A^{(S)}$ is a finite sum
$\sum_{j=1}^r a_j[s_j]$. After a covering, every $a_j$ and $s_j$ comes from
some index of the two local systems. Local filteredness supplies one common
successor for these finitely many indices, so the section is locally in the
image of $A_i^{(S_i)}$. If two such finite sums have the same image, pass
first to a local successor at which all basis elements having the same image
are identified, and group equal basis elements. The resulting finitely many
coefficient equalities hold after another local successor. Thus the map of
presheaf colimits is locally surjective and locally injective. Sheafification
gives the asserted isomorphism.
\end{proof}

\begin{theorem}
\label{theorem-illusie-I-free-abelian-n-equivalence}
Let $n \geq 2$ and let $f : X \to Y$ be an $n$-equivalence of simplicial
objects of a topos $\mathcal{T}$. Assume either
\begin{enumerate}
\item $\mathcal{T}$ has enough points, or
\item $X$ and $Y$ are simplicial abelian group objects and $f$ is a
homomorphism.
\end{enumerate}
Then the levelwise free abelian map
$$
\mathbf{Z}^{(f)} : \mathbf{Z}^{(X)} \longrightarrow \mathbf{Z}^{(Y)}
$$
is an $n$-equivalence.
\end{theorem}

\begin{proof}
Assume first that $\mathcal{T}$ has enough points. Inverse image at a point
commutes with homotopy sheaves by Lemma
\ref{lemma-illusie-I-homotopy-sheaves-colimits-pullback}, and it commutes
with free abelian objects because both are defined by colimits. Thus the
assertion at every point is Lemma
\ref{lemma-illusie-I-relative-hurewicz-free-abelian}. Points detect
epimorphisms and isomorphisms, so the result follows.

Assume now that $f$ is a map of simplicial abelian group objects. Use the
normalized complex with cohomological indexing
$N(X)^{-q} = N(X)_q$. By the Moore description of homotopy groups,
$\operatorname{Cone}(N(f))$ is acyclic in degrees at least $-n$.
Apply Lemma \ref{lemma-illusie-I-local-finite-type-approximation} with
$A = \mathbf{Z}$, $a = 0$, and $b = -n$. We obtain a locally filtered
system of maps
$$
g_i : E_i \longrightarrow F_i
$$
between complexes of finitely generated abelian groups, each with cone
acyclic in degrees at least $-n$, whose local colimit is $N(f)$.

Apply the inverse Dold--Kan functor degreewise. It uses finite direct sums,
while normalization uses finite kernels; hence both commute with the local
colimit by Lemma \ref{lemma-illusie-I-local-filtered-colimit}. We obtain
$f = \colim_L h_i$, where every $h_i = K(g_i)$ is an $n$-equivalence of
ordinary simplicial abelian groups. Lemma
\ref{lemma-illusie-I-relative-hurewicz-free-abelian} says that every
$\mathbf{Z}[h_i]$ is an $n$-equivalence. Lemma
\ref{lemma-illusie-I-local-colimit-free-modules}, applied to the constant
ring system $\mathbf{Z}$, identifies the free abelian object on the local
colimit. Homotopy sheaves commute with the local colimit by Lemma
\ref{lemma-illusie-I-local-filtered-colimit}. Therefore
$$
\mathbf{Z}^{(f)} = \colim_L \mathbf{Z}[h_i]
$$
is an $n$-equivalence.
\end{proof}

\begin{corollary}
\label{corollary-illusie-I-free-abelian-equivalence}
Under either hypothesis of Theorem
\ref{theorem-illusie-I-free-abelian-n-equivalence}, if $f$ is an
equivalence, then $\mathbf{Z}^{(f)}$ is an equivalence.
\end{corollary}

\begin{proof}
Apply the theorem for every $n \geq 2$. The conclusions for $\pi_0$ and
$\pi_1$ are already included in any one of these applications.
\end{proof}
